
que nous noterons
\[
 (x,y)\rightsquigarrow \langle x,y\rangle\in\mathbb Z
 \qquad (i\in\mathbb Z,\ x\in\Gr^i(X),\ y\in\Gr_i(X)).
\]

\textbf{Exemple 4.2.} Supposons que le schéma \(X\) soit de dimension \(n\). Prenons un élément \(y\in\Gr_n(X)\), par exemple \(y=\cl_\bullet(\underline{O}_Z)\), ou \(Z\) est un sous-schéma intègre de dimension \(n\) de \(X\). Prenons aussi un polynôme homogène de poids \(n\)
\[
 P\in\mathbb Z[T_1,\ldots,T_n],
\]
où \(T_i\) est muni du poids \(i\) \((i=1,\ldots,n)\), par exemple un monôme \(T_{i_1}\cdot T_{i_2}\cdots T_{i_k}\) avec \(i_1+i_2+\cdots+i_k=n\). Lorsqu'on substitue dans \(P\) pour chaque variable \(T_i\) \((i=1,\ldots,n)\) la classe de Chern
\[
 c_i(x)=\cl^{\bullet}(\gamma^i(x-\epsilon(x))\in\Gr^i(X)
 \qquad(x\in K^{\bullet}(X)),
\]
on obtient un élément
\[
 P(c_1(x),\ldots,c_n(X))\in\Gr^n(X).
\]
On appelle \emph{nombre de Chern de} \(x\) \emph{relativement à} \(y\) l'entier
\[
 P_y(x)=\bigl\langle P(c_1(x)),\ldots,c_n(x),y\bigr\rangle.
\]
Lorsque \(y=\cl_\bullet(\underline{O}_X)\), on omet la mention de \(y\) dans la terminologie précédente.

\subsection*{4.3.}
Sans hypothèse sur la dimension de \(X\), considérons une suite \(\underline{L}_1,\ldots,\underline{L}_d\) de \(\underline{O}_X\)-modules inversibles et un \(\underline{O}_X\)-module cohérent \(\underline{F}\) tel que \(\dim(\Supp(\underline{F}))\leq d\). Posons \(L_i=\cl^{\bullet}(\underline L_i)\) \((i=1,\ldots,d)\) et \(F=\cl_\bullet(\underline F)\).

